\documentclass[11pt,a4paper]{article}

\usepackage[utf8]{inputenc}
\usepackage[T1]{fontenc}
\usepackage{amsmath,amssymb,amsthm}
\usepackage{geometry}
\usepackage{booktabs}
\usepackage{xcolor}
\usepackage{physics}
\usepackage{hyperref}
\usepackage{enumitem}

\geometry{margin=1in}

\newtheorem{theorem}{Theorem}[section]
\newtheorem{proposition}[theorem]{Proposition}
\newtheorem{definition}[theorem]{Definition}
\newtheorem{remark}[theorem]{Remark}

%------------------------------------------------------------------------------
% CUSTOM SYMBOL: Hb ligature approximation
% The right leg of H shares with the ascender of b
%------------------------------------------------------------------------------
\newcommand{\Hb}{\textnormal{H\kern-0.12em b}}
\newcommand{\Hbn}{\Hb_n}
\newcommand{\Hbfour}{\Hb_4}

% Old notation (for comparison)
\newcommand{\Tc}{\mathcal{T}}

% Other notation
\newcommand{\dgeo}{\delta}
\newcommand{\murad}{\mu_0}
\newcommand{\Mpl}{M_{\text{Pl}}}
\newcommand{\lp}{\ell_P}

\title{\textbf{Technical Note: Redefining the Transdimensional Constant}\\[0.5em]
\Large From $\Tc$ to Holographic Bit Density $\Hb_n$}

\author{M.\ Hubka}
\date{February 2026}

\begin{document}

\maketitle

\begin{abstract}
We redefine the ``transdimensional constant'' $\Tc = 1/(4\ln 2)$ as the \textbf{holographic bit density} $\Hbfour$ specific to 4-dimensional spacetime. The generalized formula $\Hbn = 1/(n\ln 2)$ represents the information capacity per Planck area in $n$ dimensions. This reinterpretation removes conceptual ambiguity, grounds the constant in fundamental holographic physics, and opens the framework to multi-dimensional analysis. All numerical predictions remain unchanged.
\end{abstract}

%==============================================================================
\section{The Problem with ``Transdimensional''}
\label{sec:problem}
%==============================================================================

In Dimensional Collapse Theory (DCT), the constant
\begin{equation}
\Tc = \frac{1}{4\ln 2} \approx 0.3607
\end{equation}
was introduced as a ``transdimensional constant'' governing the $D \to D-2$ collapse of spacetime at black hole horizons.

However, when $\Tc$ appeared in Quantum Spectrodynamics (QSD) as:
\begin{itemize}[nosep]
    \item The Heisenberg coupling: $J = \Tc$
    \item The master relation: $\alpha^{-1} \cdot \dgeo \approx 2$ (via $\Tc/\dgeo$)
    \item The dimensionless VEV: $\nu_0 = \Tc/\dgeo$
\end{itemize}
the ``transdimensional'' interpretation became strained. QSD operates entirely within 4D; there is no dimensional collapse in the Standard Model derivations.

\textbf{The issue}: $\Tc$ is not fundamentally about dimensional \emph{transition}---it is about dimensional \emph{identity}. The ``4'' in $1/(4\ln 2)$ reflects that we live in 4D spacetime, not that dimensions are collapsing.

%==============================================================================
\section{The Holographic Bit Density}
\label{sec:definition}
%==============================================================================

\begin{definition}[Holographic Bit Density]
The \textbf{holographic bit density} in $n$-dimensional spacetime is:
\begin{equation}
\boxed{\Hbn = \frac{1}{n \ln 2}}
\end{equation}
This represents the information capacity (in bits) per Planck-area cell of an $(n-2)$-dimensional holographic screen embedded in $n$ dimensions.
\end{definition}

\begin{remark}[Notation]
We write $\Hb$ as a ligature where the right leg of `H' is shared with the ascender of `b', symbolizing the unity of \textbf{H}olographic and \textbf{b}it concepts. When a subscript is needed, we write $\Hbn$ for general dimension or $\Hbfour$ for our 4D universe.
\end{remark}

\subsection{Values Across Dimensions}

\begin{center}
\begin{tabular}{ccc}
\toprule
\textbf{Dimension $n$} & \textbf{Holographic Bit Density $\Hbn$} & \textbf{Numerical Value} \\
\midrule
2 & $1/(2\ln 2)$ & 0.7213 \\
3 & $1/(3\ln 2)$ & 0.4809 \\
\textbf{4} & $\mathbf{1/(4\ln 2)}$ & \textbf{0.3607} \\
5 & $1/(5\ln 2)$ & 0.2885 \\
6 & $1/(6\ln 2)$ & 0.2404 \\
\bottomrule
\end{tabular}
\end{center}

\textbf{Physical interpretation}: Lower-dimensional spaces have higher holographic bit density. Information is ``packed tighter'' when there are fewer dimensions to spread across.

%==============================================================================
\section{Connection to Bekenstein-Hawking Entropy}
\label{sec:BH}
%==============================================================================

The Bekenstein-Hawking entropy formula in 4D is:
\begin{equation}
S_{\text{BH}} = \frac{A}{4\lp^2} \quad \text{(in nats)}
\end{equation}

Converting to bits ($1~\text{bit} = \ln 2~\text{nats}$):
\begin{equation}
S_{\text{BH}} = \frac{A}{4\ln 2 \cdot \lp^2} \quad \text{(in bits)}
\end{equation}

This means the \textbf{area per bit} is:
\begin{equation}
A_{\text{bit}} = 4\ln 2 \cdot \lp^2
\end{equation}

The \textbf{bits per Planck area} is the inverse:
\begin{equation}
\Hbfour = \frac{1}{4\ln 2} \approx 0.3607~\text{bits}/\lp^2
\end{equation}

This is precisely our holographic bit density. The formula generalizes to $n$ dimensions as $\Hbn = 1/(n\ln 2)$, consistent with dimensional analysis of holographic entropy in arbitrary dimensions.

%==============================================================================
\section{Reinterpretation in QSD}
\label{sec:qsd}
%==============================================================================

\subsection{The Heisenberg Coupling}

\textbf{Old interpretation}: $J = \Tc$ (``transdimensional constant sets spin coupling'')

\textbf{New interpretation}: $J = \Hbfour$ (``information capacity of 4D holographic cell sets spin coupling'')

The Heisenberg coupling on $K_4$ is the rate at which quantum information can flow through each edge. This rate is bounded by the holographic bit density of the ambient spacetime.

\subsection{The Dimensionless VEV}

\textbf{Old formula}: $\nu_0 = \Tc/\dgeo$

\textbf{New formula}: $\nu_0 = \Hbfour/\dgeo$

The dimensionless vacuum expectation value is the ratio of:
\begin{itemize}[nosep]
    \item $\Hbfour$: Information capacity (what the vacuum can hold)
    \item $\dgeo$: Geometric strain (how much the code deviates from perfection)
\end{itemize}

\subsection{The Master Relation}

The relation $\alpha^{-1} \cdot \dgeo \approx 2$ becomes:
\begin{equation}
\alpha^{-1} \approx \frac{2}{\dgeo} = \frac{2 \cdot \Hbfour}{\Hbfour \cdot \dgeo} = \frac{2\Hbfour}{\nu_0^{-1} \cdot \Hbfour} 
\end{equation}

The fine structure constant is determined by the holographic information structure of 4D spacetime.

%==============================================================================
\section{Reinterpretation in DCT}
\label{sec:dct}
%==============================================================================

\subsection{The Ledger Entropy}

\textbf{Old interpretation}: Fraction $\Tc$ of horizon entropy resides on the ledger during collapse.

\textbf{New interpretation}: The ledger stores information at the 4D holographic bit density $\Hbfour$.

The ledger is not receiving a ``fraction'' of entropy---it is a 2D surface with information density $\Hbfour$ bits per Planck area, as required by 4D holography.

\subsection{Dimensional Collapse}

The $D \to D-2$ collapse still occurs, but the interpretation shifts:

\textbf{Old}: ``$\Tc$ governs the compression during collapse.''

\textbf{New}: ``Collapse occurs because different dimensions have different holographic capacities.''

The entropy change during $D \to D-2$ collapse involves:
\begin{equation}
\Delta S \propto \left(\Hb_{D-2} - \Hb_D\right) \times A = \left(\frac{1}{D-2} - \frac{1}{D}\right) \frac{A}{\ln 2}
\end{equation}

For $D=4 \to D=2$:
\begin{equation}
\Hb_2 - \Hbfour = \frac{1}{2\ln 2} - \frac{1}{4\ln 2} = \frac{1}{4\ln 2} = \Hbfour
\end{equation}

The 2D ledger has \emph{twice} the bit density of 4D spacetime, explaining why it can efficiently store the information of collapsing matter.

%==============================================================================
\section{Why This Matters}
\label{sec:significance}
%==============================================================================

\subsection{Conceptual Clarity}

\begin{center}
\begin{tabular}{lcc}
\toprule
\textbf{Aspect} & \textbf{Old ($\Tc$)} & \textbf{New ($\Hbn$)} \\
\midrule
Name & Transdimensional constant & Holographic bit density \\
Interpretation & Process (collapse) & State (structure) \\
Dimension dependence & Hidden in ``4'' & Explicit in $n$ \\
Physical meaning & Unclear & Bits per Planck area \\
QSD context & Awkward & Natural \\
\bottomrule
\end{tabular}
\end{center}

\subsection{Generalization to Other Dimensions}

The framework now naturally extends to other dimensional phases:
\begin{itemize}[nosep]
    \item \textbf{Pre-collapse phase} ($D > 4$): Lower bit density $\Hb_D < \Hbfour$
    \item \textbf{Ledger phase} ($D = 2$): Higher bit density $\Hb_2 = 2\Hbfour$
    \item \textbf{Hypothetical 3D phase}: Intermediate $\Hb_3 = (4/3)\Hbfour$
\end{itemize}

\subsection{Connection to Other Theories}

The holographic bit density connects naturally to:
\begin{itemize}[nosep]
    \item \textbf{AdS/CFT}: Boundary information density
    \item \textbf{Loop Quantum Gravity}: Spin network state counting
    \item \textbf{String Theory}: Worldsheet information capacity
    \item \textbf{Causal Set Theory}: Sprinkling density
\end{itemize}

%==============================================================================
\section{Summary of Changes}
\label{sec:summary}
%==============================================================================

\subsection{What Changes}

\begin{enumerate}
    \item \textbf{Terminology}: ``Transdimensional constant $\Tc$'' → ``Holographic bit density $\Hbn$''
    
    \item \textbf{Symbol}: $\Tc$ → $\Hbfour$ (for 4D) or $\Hbn$ (general)
    
    \item \textbf{Interpretation}: Process constant → State parameter
    
    \item \textbf{Formula}: Fixed $1/(4\ln 2)$ → General $1/(n\ln 2)$
\end{enumerate}

\subsection{What Does Not Change}

\begin{enumerate}
    \item \textbf{Numerical value}: $\Hbfour = 1/(4\ln 2) = 0.3607$ (unchanged)
    
    \item \textbf{All QSD predictions}: $\alpha^{-1}$, $\sin^2\theta_W$, $m_H$, $E_{\text{syn}}$ (unchanged)
    
    \item \textbf{All DCT results}: Ledger dynamics, UHP, entropy relations (unchanged)
    
    \item \textbf{Derived quantities}: $\dgeo$, $\murad$, $\nu_0$ (unchanged)
\end{enumerate}

%==============================================================================
\section{Updated Notation Guide}
\label{sec:notation}
%==============================================================================

\begin{center}
\begin{tabular}{lll}
\toprule
\textbf{Symbol} & \textbf{Definition} & \textbf{Value (4D)} \\
\midrule
$\Hbn$ & $1/(n\ln 2)$ & --- \\
$\Hbfour$ & $1/(4\ln 2)$ & 0.3607 \\
$\dgeo$ & $3/(80\varphi^2)$ & 0.01432 \\
$\nu_0$ & $\Hbfour/\dgeo$ & 25.18 \\
$\murad$ & $\Mpl/2^{60} \cdot \varphi^2(1-6\dgeo)$ & 25.34 GeV \\
\bottomrule
\end{tabular}
\end{center}

\subsection{LaTeX Implementation}

For documents using this notation:
\begin{verbatim}
% Hb ligature (approximate)
\newcommand{\Hb}{\textnormal{H\kern-0.12em b}}
\newcommand{\Hbn}{\Hb_n}
\newcommand{\Hbfour}{\Hb_4}
\end{verbatim}

%==============================================================================
\section{Conclusion}
\label{sec:conclusion}
%==============================================================================

The redefinition from ``transdimensional constant'' to ``holographic bit density'' is not merely cosmetic. It:

\begin{enumerate}
    \item \textbf{Grounds the constant in physics}: $\Hbfour$ has a clear meaning (bits per Planck area in 4D)
    
    \item \textbf{Removes conceptual strain}: No more invoking ``dimensional collapse'' in static QSD formulas
    
    \item \textbf{Enables generalization}: The framework naturally extends to other dimensions via $\Hbn = 1/(n\ln 2)$
    
    \item \textbf{Preserves all results}: Every number, prediction, and derivation remains valid
\end{enumerate}

\vspace{1em}
\hrule
\vspace{1em}

\textbf{Recommendation}: All future documents should use $\Hbfour$ (or $\Hbn$ in general contexts) instead of $\Tc$. The term ``transdimensional constant'' should be retired.

\end{document}
